%&template
\documentclass{article}
\usepackage{src/template}

\endofdump
\usepackage{minted}

\date{2025 -- 12 -- 07}
\useTemplate[NPA, Exercise=7]
\title{NPA -- Exercise 7}


\begin{document}
\maketitle
    
    \section{}
\begin{verbatim}
#!/usr/bin/env bash
# Group Q: 
#  Emily Seebeck - 411364
#  Michel Schnitker - 410230
#  Jarne Bärmann - 409434
#  Raphael Krafczyk - 393325

# exercise a)
for ns in R1 R2 R3 R4 R5 H1 H2; do
    ip netns del $ns 2>/dev/null
done

for ns in R1 R2 R3 R4 R5 H1 H2; do
    ip netns add $ns
done

# DEBUG: list namespaces so we know they were created

# ip netns list
# echo ""

# exercise b)
ip link add vH1-R1 type veth peer name vR1-H1
ip link add vH2-R5 type veth peer name vR5-H2
ip link add vR1-R2 type veth peer name vR2-R1
ip link add vR1-R4 type veth peer name vR4-R1
ip link add vR2-R3 type veth peer name vR3-R2
ip link add vR3-R5 type veth peer name vR5-R3
ip link add vR5-R4 type veth peer name vR4-R5

ip link set vH1-R1 netns H1
ip link set vH2-R5 netns H2
ip link set vR1-R2 netns R1
ip link set vR1-R4 netns R1
ip link set vR1-H1 netns R1
ip link set vR2-R1 netns R2
ip link set vR2-R3 netns R2
ip link set vR3-R5 netns R3
ip link set vR3-R2 netns R3
ip link set vR4-R5 netns R4
ip link set vR4-R1 netns R4
ip link set vR5-R4 netns R5
ip link set vR5-H2 netns R5
ip link set vR5-R3 netns R5

# DEBUG: Output for b)

# list links for each NS
# for NAMESPACE in R1 R2 R3 R4 R5 H1 H2; do
#     echo $NAMESPACE
#     sudo ip netns exec $NAMESPACE ip -br a
#     echo ""
# done


# exercise c) Configure ips
# Convention: IP ending with .20X for each Router RX
# Router loopbacks: fc00:X::1/128 where X is the router number
# Host IP ending with .X for HX

# H1
ip netns exec H1 ip addr add 2001:db8:1::1/64 dev vH1-R1
ip netns exec H1 ip link set dev vH1-R1 up
ip netns exec H1 ip link set lo up

# R1
ip netns exec R1 ip addr add 2001:db8:1::201/64 dev vR1-H1
ip netns exec R1 ip addr add 2001:db8:2::201/64 dev vR1-R2
ip netns exec R1 ip addr add 2001:db8:4::201/64 dev vR1-R4


ip netns exec R1 ip link set dev vR1-H1 up
ip netns exec R1 ip link set dev vR1-R2 up
ip netns exec R1 ip link set dev vR1-R4 up
ip netns exec R1 ip link set dev lo up

# R2
ip netns exec R2 ip addr add 2001:db8:2::202/64 dev vR2-R1
ip netns exec R2 ip addr add 2001:db8:3::202/64 dev vR2-R3


ip netns exec R2 ip link set dev vR2-R1 up
ip netns exec R2 ip link set dev vR2-R3 up
ip netns exec R2 ip link set dev lo up

# R3
ip netns exec R3 ip addr add 2001:db8:3::203/64 dev vR3-R2
ip netns exec R3 ip addr add 2001:db8:5::203/64 dev vR3-R5


ip netns exec R3 ip link set dev vR3-R2 up
ip netns exec R3 ip link set dev vR3-R5 up
ip netns exec R3 ip link set dev lo up

# R5
ip netns exec R5 ip addr add 2001:db8:7::205/64 dev vR5-H2
ip netns exec R5 ip addr add 2001:db8:5::205/64 dev vR5-R3
ip netns exec R5 ip addr add 2001:db8:6::205/64 dev vR5-R4


ip netns exec R5 ip link set dev vR5-H2 up
ip netns exec R5 ip link set dev vR5-R3 up
ip netns exec R5 ip link set dev vR5-R4 up
ip netns exec R5 ip link set dev lo up

# R4
ip netns exec R4 ip addr add 2001:db8:4::204/64 dev vR4-R1
ip netns exec R4 ip addr add 2001:db8:6::204/64 dev vR4-R5


ip netns exec R4 ip link set dev vR4-R1 up
ip netns exec R4 ip link set dev vR4-R5 up
ip netns exec R4 ip link set dev lo up

# H2
ip netns exec H2 ip addr add 2001:db8:7::2/64 dev vH2-R5
ip netns exec H2 ip link set dev vH2-R5 up
ip netns exec H2 ip link set lo up


# DEBUG: Output for c)
# for NAMESPACE in R1 R2 R3 R4 R5 H1 H2; do
#     echo $NAMESPACE
#     sudo ip netns exec $NAMESPACE ip a
#     echo ""
# done

# exercise d) Default Routes
ip netns exec H1 ip -6 route add default via 2001:db8:1::201
ip netns exec H2 ip -6 route add default via 2001:db8:7::205

# exercise e) Router Loopback Addresses
ip netns exec R1 ip addr add fc00:1::1/128 dev lo
ip netns exec R2 ip addr add fc00:2::1/128 dev lo
ip netns exec R3 ip addr add fc00:3::1/128 dev lo
ip netns exec R5 ip addr add fc00:5::1/128 dev lo
ip netns exec R4 ip addr add fc00:4::1/128 dev lo


# exercise f) define a route for each sid

# R1
ip netns exec R1 ip -6 route add fc00:2::1/128 via 2001:db8:2::202
ip netns exec R1 ip -6 route add fc00:4::1/128 via 2001:db8:4::204

# R2
ip netns exec R2 ip -6 route add fc00:1::1/128 via 2001:db8:2::201
ip netns exec R2 ip -6 route add fc00:3::1/128 via 2001:db8:3::203

# R3
ip netns exec R3 ip -6 route add fc00:2::1/128 via 2001:db8:3::202
ip netns exec R3 ip -6 route add fc00:5::1/128 via 2001:db8:5::205

# R5
ip netns exec R5 ip -6 route add fc00:3::1/128 via 2001:db8:5::203
ip netns exec R5 ip -6 route add fc00:4::1/128 via 2001:db8:6::204

# R4
ip netns exec R4 ip -6 route add fc00:1::1/128 via 2001:db8:4::201
ip netns exec R4 ip -6 route add fc00:5::1/128 via 2001:db8:6::205


# exercise g) enable forwarding on router

sysctl -w net.ipv6.conf.all.seg6_enabled=1
sysctl -w net.ipv6.conf.default.seg6_enabled=1

for ROUTER in R1 R2 R3 R4 R5; do
    ip netns exec $ROUTER sysctl -w net.ipv6.conf.all.forwarding=1
    ip netns exec $ROUTER sysctl -w net.ipv6.conf.all.seg6_enabled=1
done

ip netns exec R1 sysctl -w net.ipv6.conf.vR1-H1.seg6_enabled=1
ip netns exec R1 sysctl -w net.ipv6.conf.vR1-R2.seg6_enabled=1
ip netns exec R1 sysctl -w net.ipv6.conf.vR1-R4.seg6_enabled=1

ip netns exec R2 sysctl -w net.ipv6.conf.vR2-R1.seg6_enabled=1
ip netns exec R2 sysctl -w net.ipv6.conf.vR2-R3.seg6_enabled=1

ip netns exec R3 sysctl -w net.ipv6.conf.vR3-R2.seg6_enabled=1
ip netns exec R3 sysctl -w net.ipv6.conf.vR3-R5.seg6_enabled=1

ip netns exec R4 sysctl -w net.ipv6.conf.vR4-R1.seg6_enabled=1
ip netns exec R4 sysctl -w net.ipv6.conf.vR4-R5.seg6_enabled=1

ip netns exec R5 sysctl -w net.ipv6.conf.vR5-H2.seg6_enabled=1
ip netns exec R5 sysctl -w net.ipv6.conf.vR5-R3.seg6_enabled=1
ip netns exec R5 sysctl -w net.ipv6.conf.vR5-R4.seg6_enabled=1

# DEBUG: Output for g)
# for ROUTER in R1 R2 R3 R4 R5; do
#     ip netns exec $ROUTER ip a
# done

# exercise h) End Action
ip netns exec R1 ip route add fc00:1::1/128 encap seg6local action End dev lo
ip netns exec R2 ip route add fc00:2::1/128 encap seg6local action End dev lo
ip netns exec R3 ip route add fc00:3::1/128 encap seg6local action End dev lo
ip netns exec R4 ip route add fc00:4::1/128 encap seg6local action End dev lo
ip netns exec R5 ip route add fc00:5::1/128 encap seg6local action End dev lo


# exercise i) Routing
ip netns exec R1 ip -6 route add 2001:db8:7::/64 encap seg6 mode encap segs fc00:2::1,fc00:3::1,fc00:5::1 dev vR1-R2
ip netns exec R5 ip -6 route add 2001:db8:1::/64 encap seg6 mode encap segs fc00:4::1,fc00:1::1 dev vR5-R4

# exercise j) Ping Test

ip netns exec H1 ping6 -c 5 2001:db8:7::2

# PING 2001:db8:7::2(2001:db8:7::2) 56 data bytes
# 64 bytes from 2001:db8:7::2: icmp_seq=1 ttl=63 time=0.101 ms
# 64 bytes from 2001:db8:7::2: icmp_seq=2 ttl=63 time=0.146 ms
# 64 bytes from 2001:db8:7::2: icmp_seq=3 ttl=63 time=0.131 ms
# 64 bytes from 2001:db8:7::2: icmp_seq=4 ttl=63 time=0.124 ms
# 64 bytes from 2001:db8:7::2: icmp_seq=5 ttl=63 time=0.126 ms

# --- 2001:db8:7::2 ping statistics ---
# 5 packets transmitted, 5 received, 0% packet loss, time 4073ms
# rtt min/avg/max/mdev = 0.101/0.125/0.146/0.014 ms

# exercise k) capture traffic

echo -e "\n\nStarting tcpdump..."

for ROUTER in R1 R2 R3 R4 R5; do
    ip netns exec $ROUTER ip -o link show | awk -F': ' '{print $2}' | sed 's/@.*$//' | grep -v '^lo$' | while read -r IF; do
        ip netns exec $ROUTER tcpdump -v -n -w $ROUTER_$IF.pcap -i $IF &
    done
done

sleep 1

echo -e "\n\nStarting ping..."

NUM_PINGS=5
ip netns exec H1 ping6 -c $NUM_PINGS 2001:db8:7::2

sleep $(($NUM_PINGS+3))
killall tcpdump



# exercise l) tshark



verify_requests() {
    # $1 = filename
    if [ $(tshark -r "$1" -Y "icmpv6.type == 128" | wc -l) == 0 ]; then
        echo "tshark output in $1 contains no icmp requests where some should have been"
        exit 1
    fi
}

verify_no_requests() {
    if [ $(tshark -r "$1" -Y "icmpv6.type == 128" | wc -l) != 0 ]; then
        echo "tshark output in $1 contains icmp requests where none should have been"
        exit 1
    fi
}

verify_responses() {
    if [ $(tshark -r "$1" -Y "icmpv6.type == 129" | wc -l) == 0 ]; then
        echo "tshark output in $1 contains no icmp responses where some should have been"
        exit 1
    fi
}

verify_no_responses() {
    if [ $(tshark -r "$1" -Y "icmpv6.type == 129" | wc -l) != 0 ]; then
        echo "tshark output in $1 contains icmp responses where none should have been"
        exit 1
    fi
}

# Request: H1 → H2
verify_requests "vR1-H1.pcap"
verify_responses "vR1-H1.pcap"

verify_requests "vR1-R2.pcap"
verify_no_responses "vR1-R2.pcap"

verify_requests "vR2-R3.pcap"
verify_no_responses "vR2-R3.pcap"

verify_requests "vR3-R5.pcap"
verify_no_responses "vR3-R5.pcap"

verify_requests "vR5-H2.pcap"
verify_responses "vR5-H2.pcap"

# Reply: H2 → H1
verify_no_requests "vR5-R4.pcap"
verify_responses "vR5-R4.pcap"

verify_no_requests "vR4-R1.pcap"
verify_responses "vR4-R1.pcap"

echo "OK"


\end{verbatim}


\end{document}